\scanpage{3}
\mspage{2}

\deleted{De plus, on veut que la structure implique}

\deleted{\textcircled{\scriptsize b}\quad $\forall\,T\in P_{f\,\mathrm{mod}}(\mathcal L)$, on se donne une
relation d'équivalence $R_T$ dans l'ens
\unsure{où}}

Si $T$ est une partie finie
modérée\unsure{,|;}, on \deleted{\unsure{désignera|désigne} par} \inserted{pose}

\[
 C\deleted{\unsure{\mathrm{omp}^{r}|\mathrm{omp}}}_{\mathcal L}(T)=
 \{\,t\in\mathcal L\setminus T\mid \{T\cup\{t\}\}\ \text{\unsure{modéré|modérée}}\,\}
\]

% slanted
\margin{\uline{le complémentaire} \uline{modéré} de $T$.
Question : $T$ est-il connu quand on \unsure{connait|connaît} $C_{\mathcal L}(T)$ ? Plus précisément,
si $T'\subset T$, on a $C(T)\subset C(T')$, l'inverse est-il vrai ? Oui,
Donc, si \inserted{$t\in T'$} $t\notin T$, \deleted{\unsure{$t$ est gén.} \unsure{tel} ($T\cup\{t\}$ modéré) \unsure{puisque} \unsure{on a toujours} \unsure{$s\notin T$,}} alors $t\in C(T)$ et $t\notin C(T')$}

ou mieux, l'ens des $t\in L\setminus T$ en
position modérée p.r. à $T$.

[\uline{NB}\quad deux parties finies modérées
sont dites en position relative
modérée si leur-réunion est
modérée ; $t$ est dit en pos. rel.
modérée p.r. à $T$ si $\{t\}$ l'est ;
$s$ et $t$ sont dits en pos. rel. modérée
si $\{s,t\}$ est modéré i.e. \deleted{\unsure{si}},
$s\ne t$ ou $\{s;t\}\in S_0$ \ldots\ ]

% slanted
\margin{\uline{NB} $T$\deleted{\unsure{est}} est connu quand on \unsure{connait|connaît} $C(T)$}

\textcircled{\scriptsize b}\quad $\forall\,T\in\underbrace{\mathrm{Modf}}_{\substack{\text{parties mod}\\\text{finies}}}(\mathcal L)$, une relation
d'équivalence $R_T$ dans $C_{\unsure{\checkmark|\mathcal L|}}(T\unsure{)|\}}$ (dont
les classes s'appellent les \uline{comp. connexes} \inserted{de $C(T)$}

% paper watermark, visible from the reverse
\rotatebox{90}{\textsc{\unsure{xtra Strong Ext}}}
